% Source: 04 -Mon Oct 5 2026.pdf, page 10. Content remains in source order.
\sourcepage{04}{10}
An algorithm executing $t$ steps can be simulated by cascading $t$ copies of $C_{\mathrm{RAM}}$:
\sourcefigure[0.95]{assets/04/page-010-cascaded-ram.png}

\subsection{Crucial points:}
\begin{enumerate}
\item $|\langle C_{\mathrm{RAM}}\rangle|$ is polynomial in $x$ and $y$.
\begin{itemize}\item circuitry polynomial in the memory size \textbf{NON CHIARO}.\end{itemize}
\item It suffices to set $|y|=\Theta(|x|^{k_1})$. 
\item(we can avoid examining every possible input length : we can limit from above the quantity of input wires.)
\item the number of wires is polynomial : parallelism with the implementation of the ALU. The Sum component which is  the basic component to work needs a number of bits that is at most quadratic in the size of the operands. 
\item It suffices to consider $t=\Theta((|x|+|y|)^{k_2})$. 
\item (We will not need more than $(|x|+|y|)^{k_2}$ copies of the circuit to run the verification algorithm, because we know that since L is in NP it will run in polynomial time )
\end{enumerate}
($k_1,k_2$ from definition of poly-time verifier.)

We have a polynomial number of circuits with a polynomial number of wires : the product of a polynomial is a polynomial $\Rightarrow C_{V_L,x}$ has size polynomial in $x$!
Therefore $f_{L\to\mathrm{BC\text{-}SAT}}$ is PTC.
